\section{REVIEW OF LITERATURE}

\subsection{Code-Switching and Taglish}

Code-switching is a common feature of multilingual communities.
It has been characterized as a systematic linguistic practice
rather than a random mixture of languages \cite{b2}. In the
Philippine setting, Tagalog-English code-switching, commonly
referred to as Taglish, has been described as a meaningful mode
of discourse that reflects social context, identity, and
communicative purpose \cite{b1,b3}.

For NLP systems, Taglish is difficult because it combines
tokens, grammatical structures, and discourse cues from more
than one language. A sentence may contain English content
words, Filipino function words, localized slang, and
non-standard spelling in a single utterance. This makes it
challenging for tools trained only on formal Filipino or formal
English to correctly interpret meaning.

\subsection{Informal Language and Internet Slang}

Internet communication encourages rapid lexical innovation.
Digital environments have been shown to create new forms of
written expression because users adapt language to speed,
identity, humor, and platform conventions \cite{b4}. In
Filipino online spaces, these conditions produce expressions
that may be short-lived, community-specific, or highly
dependent on current events.

Slang is especially difficult to model because its meaning is
not always recoverable from the word form alone. Some terms are
structurally novel, while others are existing words with
shifted meanings. For instance, an English word may be borrowed
into Filipino online discourse and acquire a new pragmatic
meaning. Therefore, detecting slang requires more than checking
whether a token appears in a dictionary.

\subsection{Low-Resource Filipino NLP}

Filipino NLP remains a relatively low-resource area compared
with English and other widely studied languages. The
availability of tools such as calamanCy provides an important
foundation for Tagalog processing because it offers
spaCy-based pipelines for tokenization and other NLP tasks
\cite{b6}. However, informal Filipino and Taglish still pose
additional challenges because social media text often contains
code-switching, creative spelling, abbreviations, and emerging
vocabulary.

For systems such as Pinoy Speak, low-resource NLP tools must
be combined with domain-specific heuristics and corpus-based
analysis. This is necessary because a general Filipino NLP
pipeline may identify tokens but may not automatically
determine whether a term is slang, standard vocabulary,
profanity, a named entity, or a context-dependent semantic
shift.

\subsection{Human Annotation and Validation in Filipino NLP}

Human judgment remains important in Filipino NLP tasks that
involve context-sensitive meaning. A Filipino metaphor
identification study \cite{b12} used manual data annotation
by the researcher and four native Filipino annotators. The
annotation team included students with communication and
language-related backgrounds, and the process involved
checking and normalizing Filipino text before model training.
This demonstrates that computational processing alone may be
insufficient when the target phenomenon depends on context,
figurative use, or local linguistic knowledge.

This is relevant to the present study because Filipino internet
slang is also context-sensitive. A term may be ordinary
vocabulary in one sentence but slang in another. Some words
also require cultural familiarity to determine whether the
definition, usage example, and formation type are acceptable.
Therefore, this study uses Filipino or Taglish-speaking
validators to assess the generated lexicon entries rather than
relying only on automatic detection.

\subsection{Subword-Aware Embeddings and OOV Handling}

Subword-aware embeddings are useful for handling rare,
misspelled, and morphologically varied words. FastText word
representations, which model words as bags of character
n-grams, were introduced to address this challenge \cite{b5}.
This approach allows a system to generate useful
representations even for words that are rare or unseen during
training.

This property is relevant to informal Filipino text because
slang often appears in creative or non-standard forms.
Character n-gram representations can provide nearest-neighbor
evidence for candidate terms, helping the system determine
whether a word resembles known slang, standard Filipino
morphology, or unrelated noise.

\subsection{Semantic Shift and Diachronic Language Change}

Semantic shift refers to changes in word meaning over time or
across contexts. Diachronic word embeddings have been shown to
reveal statistical patterns of semantic change \cite{b7}.
Methods for tracking semantic shifts using word embeddings have
also been surveyed \cite{b8}, and dynamic word embeddings for
modeling language change through time have been proposed
\cite{b9}.

These studies support the idea that language change can be
observed computationally through distributional patterns. For
Filipino slang detection, this is important because some terms
are not new words but old words used in new ways. A
trend-analysis system must therefore consider both novel
out-of-vocabulary words and known words that may be undergoing
contextual or semantic change.

\subsection{Retrieval-Augmented Generation}

Large Language Models can generate fluent explanations, but
they may also hallucinate or provide unsupported definitions
when asked about specialized or newly emerging terms.
Retrieval-Augmented Generation addresses this limitation by
combining generation with retrieval from an external knowledge
source \cite{b10}. Instead of relying only on model memory, a
RAG system retrieves relevant documents or entries and uses
them as context for the generated response.

This is useful for a slang tutor because the chatbot can
ground its explanations in the live lexicon. In Pinoy Speak,
the tutor module retrieves relevant dictionary entries before
producing responses, reducing the risk of inventing
unsupported slang meanings.

\subsection{Usability Evaluation}

System usability is an important part of evaluating applied
software systems. Brooke's System Usability Scale provides a
widely used framework for measuring perceived usability through
user responses \cite{b11}. Although this study does not
implement the full SUS scoring procedure, it follows the same
principle of collecting user ratings on navigation, clarity,
usefulness, and confidence after system use. This supports the
evaluation of Pinoy Speak not only as an NLP pipeline but also
as a usable educational web application.
